% =============================================================================
% Kapitel 4: RFC-Analyse und Anforderungsmodell
% =============================================================================

\chapter{Analyse und Anforderungsmodell}
\label{chap:analyse}

\chapref{chap:methodik} legt die Quellenregeln fest. Dieses Kapitel wendet sie auf RFC~9868 an
(\secref{sec:extraktionsmethode}), überführt den Anforderungskatalog in den FF1-Maßstab
(\secref{sec:anforderungsmatrix}) und bildet das FF2-Pfad- und Auswertungsmodell
(\secref{sec:ff2-pfadmodell}).
Die Technologieauswahl (\secref{sec:technologieauswahl}) schließt die Analyse ab.

\section{Anwendung der Extraktionsmethode}
\label{sec:extraktionsmethode}

Die Auswertung wendet die in \secref{sec:rfc-auswertung} festgelegte Methode auf den Primärtext an. Die
103 mit \texttt{>{}>} markierten Absätze \cite[Sec.~2]{rfc9868} dienen als Kontrollzahl; unmarkierte
normative Festlegungen bleiben Teil des Maßstabs \cite{rfc2119,rfc8174}. Die Auswertung erfasst Aussagen
Satz für Satz statt Absatz für Absatz, denn ein markierter Absatz kann mehrere Pflichten mit verschiedenen
Adressaten bündeln. Jede erfasste Aussage wird zu einem Prüfpunkt für den Anforderungskatalog der
Implementierung (\secref{sec:anforderungsmatrix}).

\figref{fig:normverteilung} zeigt, wie sich die 103 markierten Absätze auf die Abschnitte verteilen:

\begin{figure}[H]
  \centering
  \begin{tikzpicture}[x=1cm, y=1cm]
    \draw[sqdunkel, thin] (-0.35,0) -- (13.35,0);
    \foreach \s/\v [count=\i from 0] in {1/0,2/0,3/0,4/0,5/0,6/0,7/0,8/2,9/5,10/21,11/44,12/3,13/9,14/6,%
      15/5,16/1,17/0,18/0,19/2,20/0,21/0,22/0,23/0,24/0,25/4,26/1,27/0}{
      \node[font=\tiny, text=feldlinie] at (\i*0.5,-0.25) {\s};
      \ifnum\v>0
        \draw[draw=sqblaurand, fill=sqblau] (\i*0.5-0.17,0) rectangle (\i*0.5+0.17,\v*0.088);
        \node[font=\tiny, text=sqdunkel] at (\i*0.5,\v*0.088+0.18) {\v};
      \fi
    }
    \draw[decorate, decoration={brace, mirror, amplitude=4pt}, draw=sqgrau, thin] (-0.2,-0.5) -- (3.2,-0.5);
    \node[font=\scriptsize, text=feldlinie, align=center] at (1.5,-0.95) {Abschnitte 1 bis 7:\\ keine};
    \node[font=\scriptsize, text=feldlinie] at (8.85,-0.85) {Nummer des RFC-Abschnitts};
  \end{tikzpicture}
  \caption{Verteilung der markierten normativen Absätze über die Abschnitte}
  \label{fig:normverteilung}
\end{figure}

Knapp zwei Drittel der markierten Absätze (65 von 103) liegen in den Abschnitten~10 und 11. Sie betreffen
das Optionsformat, seine Verarbeitung und die einzelnen Optionen. Bis einschließlich Abschnitt~7 gibt
es keine markierten Absätze. Abschnitt~7 enthält dennoch eine für die Arbeit relevante Längenregel:
Zusätzliche Header zwischen IPv4- und UDP-Header, sogenannte Shim-Header, verringern die Obergrenze der
\texttt{UDP Length} um ihre Gesamtlänge (\secref{sec:ipv4}). Die Implementierung verarbeitet nur
\texttt{Protocol}~17 und verfolgt keine Shim-Ketten. Das ist eine Umfangsgrenze dieser Arbeit.

\section{Anforderungskatalog und FF1-Maßstab}
\label{sec:anforderungsmatrix}
\label{sec:ff1-kriterien}

Das Repository der Implementierung (\chapref{chap:entwurf}) führt in \texttt{docs/requirements.md} 49
funktionale Anforderungen (FR) und zwölf nichtfunktionale (NFR). Die Datei entstand vor dem ersten Funktionsschritt
und wurde während der Umsetzung fortgeschrieben.

Die Konformitätsmatrix in dieser Datei umfasst 67 Zeilen. Sie dient als \textbf{RFC-Arbeitsindex}:
Sie ordnet Aussagen des RFC den Anforderungen und ihrem Um\-setzungs\-stand zu. Neben Normsätzen enthält sie
Beschreibungen und Leitlinien. Die Prüfung aus \secref{sec:extraktionsmethode} stuft deshalb jede Zeile
erneut am Primärtext ein. Normstufe, Anwendbarkeit, FF1-Bewertung und Implementierungs\-stand
(\secref{sec:soll-ist}) stehen in einer eigenen Datei, der \textbf{geprüften
Einzelzuordnung}.\footnote{\raggedright Die Datei
\nolinkurl{thesis/daten/konformitaet-kategorien.csv} liegt im öffentlichen Repository dieser Arbeit,
\url{https://github.com/ab7z/mcs-thesis-docs} (Stand der abgegebenen Fassung).}
Zeilen, die mehrere Normsätze oder Adressaten bündeln, etwa die Zeilen zu den Abschnitten 11.2, 16 und 19
sowie 25.2 bis 25.4, bleiben eine Zeile. Bewertet wird jeweils der Anteil, der die Endpunkte im Sinne
von FF1 betrifft; die Transitpflicht aus Abschnitt~16 gehört zu FF2.
\tabref{tab:anforderungsmatrix} zeigt eine eigene übersetzte Auswahl; Normstufe und Abdeckung stehen im
Wortlaut des Arbeitsindex:

\begin{table}[H]
  \centering
  \caption{Auszug aus dem RFC-Arbeitsindex der Implementierung}
  \label{tab:anforderungsmatrix}
  \footnotesize
  \setlength{\tabcolsep}{4pt}
  \renewcommand{\arraystretch}{1.0}
  \begin{tabularx}{\textwidth}{@{}L{1.3cm}YL{2.25cm}L{2.1cm}L{1.9cm}@{}}
    \toprule
    \textbf{Fund\-stelle} & \textbf{Normative Aussage} & \textbf{Norm\-stufe} & \textbf{Abdeckung} &
      \textbf{FR/NFR} \\
    \midrule
    Sec.~8 & Ausrichtungsbyte vor dem OCS ist null; sonst alle Optionen ignorieren und die Surplus Area
      still verwerfen & \texttt{MUST} & \texttt{yes} & FR-05 \\
    Sec.~9 & OCS ungleich null, wenn die UDP-Prüfsumme ungleich null ist & \texttt{MUST} & \texttt{yes} &
      FR-21 \\
    Sec.~10 & Optionen über 254~Byte nutzen das erweiterte Format; das kleinste Format ist empfohlen &
      \texttt{MUST/}\allowbreak\texttt{SHOULD} & \texttt{yes (send); local strict receive} & FR-09, FR-14 \\
    Sec.~11.2 & Höchstens sieben NOPs in Folge; übermäßige Folgen protokollieren und Ressourcen begrenzen &
      \texttt{SHOULD} & \texttt{partial} & FR-16, NFR-05 \\
    Sec.~12 & UCMP und UENC sind reserviert; UEXP ist die UNSAFE-Experimentoption & \texttt{reserved} &
      \texttt{out} & FR-39 \\
    \bottomrule
  \end{tabularx}
\end{table}

Die Abdeckung ist eine Selbstauskunft des Arbeitsindex und kann von der eigenen Bewertung abweichen.
Die Zeile zu Abschnitt~11.2 verbindet die NOP-Pflichten mit der bedingten Empfehlung aus Abschnitt~25.2,
Ressourcen zu begrenzen \cite[Sec.~25.2]{rfc9868}.

Bewertet wird gegen den Maßstab von FF1: welche Anforderungen sich unter Linux und IPv4 im Benutzerraum
\zitat{vollständig, teilweise oder nicht} erfüllen lassen (\secref{sec:zielsetzung}). Die Bewertung nutzt
drei Kategorien und eine Sonderklasse, die \tabref{tab:operationalisierung} bereits verankert.
\textbf{Vollständig} erfüllbar ist eine Anforderung, wenn sie im gewählten Rahmen ohne Abstriche umsetzbar
ist. Dieser Rahmen umfasst Linux, IPv4 und Raw Sockets im Benutzerraum sowie die Beschränkung auf
\texttt{Protocol}~17 ohne Shim-Ketten (\secref{sec:zielsetzung}, \secref{sec:extraktionsmethode}).
\textbf{Teilweise} erfüllbar ist sie, wenn nur ein Teil erreichbar bleibt und die Ursache der
Einschränkung belegt ist. \textbf{Nicht erfüllbar} ist sie, wenn das dokumentierte Verhalten des
Betriebssystems oder seiner Schnittstellen die Umsetzung im Benutzerraum grundsätzlich ausschließt.
\textbf{Nicht anwendbar}
bleibt als Sonderklasse sichtbar, wird aber nicht bewertet. Hierhin gehören Aussagen außerhalb des
Umfangs, nicht gewählte \texttt{MAY}-Funktionen und reine Entwurfsleitlinien, die keine zu bewertende
Umsetzungspflicht festlegen.

\section{FF2-Pfad- und Auswertungsmodell}
\label{sec:ff2-pfadmodell}

FF2 fragt, in welchem Umfang die Surplus Area auf den untersuchten realen Netzpfaden bis zur Gegenstelle
unverändert erhalten bleibt
(\secref{sec:zielsetzung}). Ausgewertet wird je Pfad, Richtung und Messzeitpunkt ein kontrolliertes
Szenario aus Kontroll- und Optionsverkehr, das mehrere Wiederholungen oder bei FRAG mehrere
IP-Datagramme umfassen kann (\tabref{tab:operationalisierung}). Ein gültiger Lauf beschreibt die Surplus
Area als \textbf{erhalten}, \textbf{verändert} oder \textbf{entfernt} oder das Datagramm als
\textbf{verworfen}; fehlerhafte Messungen
erhalten den Status \textbf{Lauf ungültig} und gehen nicht in die Bewertung ein
(\secref{sec:testkonzept}). Die fünf \textbf{Pfadstufen} bilden getrennte Messklassen mit
unterschiedlicher Beteiligung fremder Netze:

\begin{enumerate}
  \item \textbf{Loopback:} Sende- und Empfangspfad auf demselben Host; prüft Bibliothek und Messwerkzeuge
    ohne fremde Akteure.
  \item \textbf{Lokale virtuelle Testumgebung:} Eine Ubuntu-VM unter VMware Fusion,
    mit Linux-Netzwerk-Namensräumen und \texttt{veth}; prüft Weiterleitung, Adress\-umsetzung und
    Paketfilter unter eigener Kontrolle.
  \item \textbf{Tunnel:} WireGuard kapselt das innere Datagramm samt Surplus Area. Der äußere Pfad
    transportiert die verschlüsselten Tunnelpakete zwischen den Tunnel\-endpunkten.
  \item \textbf{Öffentliche Pfade in Europa:} Cloud-Endpunkte in Deutschland und Finnland sowie ein
    Mobilfunk-Zugangsnetz; reale Zugangs-, Transit- und Provider-Kanten.
  \item \textbf{Interkontinentale Pfade:} Endpunkte in Nordamerika und Asien-Pazifik; lange Transitketten
    und die Netzstrukturen großer Cloud-Anbieter.
\end{enumerate}

Vorversuche ohne beidseitige Mitschnitte oder Sendemanifest bleiben Piloten und werden nicht
verallgemeinert (\secref{sec:testkonzept}).

Die \textbf{Auswertung} verknüpft für jedes gerichtete Szenario die vorhandenen Sende- und Empfangsbelege.
Bei vollständigen Kampagnen sind dies das Ergebnis des Sendeaufrufs, das Sendemanifest, Sender-PCAP,
Empfänger-PCAP und Empfänger\-ausgabe. Das Sendemanifest hält je Sendevorgang unter anderem Folgenummer,
Adressen und Längen fest. Zusammen helfen die Belege, Senderfehler, Abweichungen auf dem Pfad und
Verwerfungen im Empfänger zu unterscheiden.

Zusätzliche Mitschnitte, etwa am macOS-Messpunkt aus \tabref{tab:scope-messung}, können Teile des Pfads
weiter eingrenzen. Ein Befund wird nur dem Abschnitt zwischen den Messpunkten zugeschrieben, die ihn
tatsächlich begrenzen. Dieser Abschnitt heißt \textbf{Attributionsintervall}. Liegen nur Mitschnitte
an den beiden Endpunkten vor, umfasst es den gesamten Pfad dazwischen. Ein konkretes Gerät oder dessen
interner Mechanismus lässt sich daraus nicht bestimmen.

Die aus Vorversuchen und Messungen abgeleiteten Mechanismen\-klassen stellt \chapref{chap:evaluation} mit
ihren Belegen dar.

Zwei Tabellen grenzen abschließend den Geltungsbereich ab: \tabref{tab:scope-implementierung} den
Implementierungsumfang, \tabref{tab:scope-messung} die bewusst breitere Messinfrastruktur:

\begin{table}[H]
  \centering
  \caption{Implementierungsumfang}
  \label{tab:scope-implementierung}
  \small
  \begin{tabularx}{\textwidth}{@{}L{3.4cm}YY@{}}
    \toprule
    \textbf{Merkmal} & \textbf{Festlegung} & \textbf{Außerhalb des Umfangs} \\
    \midrule
    Betriebssystem   & Linux & andere Betriebssysteme \\
    Vermittlungsschicht & IPv4 & IPv6 \\
    Ausführungsort   & Benutzerraum mit Raw Sockets & Umsetzung im Kernel \\
    Artefakt         & Bibliothek mit Mess- und Beispielprogrammen & eigenständiger Dienst;
      Protokollautomaten für REQ/RES und TIME \\
    Optionsumfang    & acht Pflichtoptionen (im RFC \textit{must-support}) sowie allgemeine Empfangsregeln
      für unbekannte SAFE- und UNSAFE-Arten & typisierte Erzeugung und Verarbeitung weiterer Arten wie TIME, AUTH, EXP,
      UCMP, UENC und UEXP \\
    \bottomrule
  \end{tabularx}
\end{table}

\begin{table}[H]
  \centering
  \caption{Messinfrastruktur}
  \label{tab:scope-messung}
  \small
  \begin{tabularx}{\textwidth}{@{}L{5.4cm}Y@{}}
    \toprule
    \textbf{Baustein} & \textbf{Rolle in der Messung} \\
    \midrule
    Linux-Endpunkte: lokale virtuelle Maschine und Cloud-Instanzen in Europa, Nordamerika
      und Asien-Pazifik & Sende- und Empfangsendpunkte der Kampagnen; Mitschnitt mit \texttt{tcpdump} \\
    Mobilfunk-Zugang über einen Hotspot & reales Zugangsnetz mit Hotspot-NAT; der Anteil eines
      möglichen CGNAT bleibt ohne zusätzlichen Messpunkt offen \\
    macOS-Messpunkt \texttt{en0} (\texttt{access\_bpf}) & zusätzlicher Beobachtungspunkt am Zugangsnetz;
      die Bibliothek wird dort nicht eingesetzt \\
    WireGuard-Tunnel & Kontrollpfad mit Kapselung (Pfadstufe~3) \\
    \bottomrule
  \end{tabularx}
\end{table}

\section{Technologieauswahl}
\label{sec:technologieauswahl}

Die Implementierung entsteht als einbindbare \textit{Bibliothek} mit Mess- und Beispiel\-programmen.
Diese Form passt zu den Entwurfsprinzipien von RFC~9868: Zustände verwaltet grundsätzlich die Anwendung
oder eine Schicht oder Bibliothek in ihrem Auftrag; von dort gehen auch Antworten auf Optionen aus. Die
Reassemblierung von Fragmenten ist die einzige, begrenzte Ausnahme beim Zustand
(\secref{subsec:surplus-prinzipien}) \cite[Sec.~6]{rfc9868}. Zu begründen sind die Sprache, der Zugang
zur Surplus Area, die Messwerkzeuge sowie Fuzzing und formale Gegenproben.

\subsection{Programmiersprache}
\label{subsec:programmiersprache}

Die Sprache muss vier Bedingungen erfüllen. Erstens braucht die Bibliothek Zugriff auf das vollständige IP-Datagramm
einschließlich der Surplus Area (\secref{sec:netzpfad}). Zweitens muss der Serialisierer das Paketformat bytegenau in
Netzwerkbyteordnung schreiben; native Speicherabbilder mit möglichem Padding werden nicht als Paketformat verwendet.
Drittens verarbeitet der Parser auch Pakete unbekannter oder bösartiger Absender. An dieser \textit{Vertrauensgrenze}
muss er fehlerhafte und gezielt konstruierte Eingaben sicher behandeln. Viertens soll die Bibliothek keine verwaltete
Laufzeitumgebung voraussetzen und ohne automatische Speicher\-bereinigung (\mbox{Garbage} Collector) auskommen. Das
erleichtert ihre Einbindung und vermeidet GC-Pausen, garantiert aber kein allgemein vorhersagbares Zeitverhalten. Die
ersten beiden Bedingungen betreffen den Protokollzugriff, die letzten beiden sind Qualitätsziele der Arbeit.

C, C++, Zig und Rust erlauben den benötigten Zugriff auf Speicher und System\-schnittstellen und
erfüllen die Vorgabe, ohne Garbage Collector auszukommen. Go und Java werden wegen dieser ausdrücklich
gewählten Vorgabe nicht weiter betrachtet.

Den Ausschlag gibt die Speicher\-sicherheit an der Vertrauensgrenze. Sie schließt Zugriffe außerhalb eines
Speicher\-bereichs (\textit{Buffer Overflow}), auf freigegebenen Speicher (\textit{Use-after-free}) und doppelte
Freigaben (\textit{Double Free}) aus. Ein \textit{Memory Leak} gewährt dagegen keinen unzulässigen Zugriff und bleibt
auch in Rust möglich. Das Microsoft Security Response Center (MSRC) berichtete 2019, dass rund 70~Prozent der
Schwachstellen, denen Microsoft jährlich eine CVE-Nummer (Common Vul\-ner\-a\-bil\-i\-ties and Exposures) zuweist, auf
Speicher\-fehler zurückgehen \cite{msrc-memory-blog}. Chromium nennt rund 70~Prozent unter 912 seit 2015 erfassten
Fehlern hoher oder kritischer Schwere in der stabilen Auslieferung (Stable Channel) \cite{chromium-memory}. Die National
Security Agency (NSA) empfiehlt speichersichere Sprachen, darunter Rust \cite{nsa-memory-safety}.

Die für die Entscheidung wichtigen Eigenschaften fasst \tabref{tab:sprachvergleich} zusammen
\cite{iso-c,cpp-core-guidelines,zig-lang,rust-book}:

\begin{table}[H]
  \centering
  \caption{Vergleich der Sprachkandidaten}
  \label{tab:sprachvergleich}
  \footnotesize
  \begin{tabularx}{\textwidth}{@{}L{2.7cm}YYYY@{}}
    \toprule
    \textbf{Kriterium} & \textbf{C} & \textbf{C++} & \textbf{Zig} & \textbf{Rust} \\
    \midrule
    Speicher\-sicherheit & nicht garantiert & nicht garantiert & nicht vollständig; Prüfungen je nach
      Build-Modus & garantiert in Safe Rust; bei \texttt{unsafe} nur unter Bedingungen \\
    Speicher\-verwaltung & manuell & manuell oder an Objekt\-lebens\-dauer gebunden & explizite Zuteilung
      und Freigabe & Ownership \\
    \bottomrule
  \end{tabularx}
\end{table}

Unter den vier Kandidaten verbindet allein Rust den Schutz vor unzulässigen Speicherzugriffen im sicheren
Sprachumfang mit einer Speicherverwaltung ohne Garbage Collector. Ein Zugriff mit einem ungültigen Array-
oder Slice-Index löst einen
definierten \textit{panic} aus \cite{rust-lang,rust-book}. In C ist ein Zugriff außerhalb der Grenzen
eines Arrays undefiniertes Verhalten \cite{iso-c}. Die Rust-Garantien setzen voraus, dass die
Sicherheitsbedingungen an jeder \texttt{unsafe}-Grenze eingehalten werden.\footnote{Das Rust-Schlüsselwort
\texttt{unsafe} ist von der
UNSAFE-Optionsklasse aus \secref{subsec:surplus-prinzipien} zu unterscheiden.}
Die Regeln zu Ownership und Lebensdauer prüft der Übersetzer statisch, also ohne den Quelltext
auszuführen \cite{rust-book}; dynamische Prüfungen und Fuzzing erfassen dagegen nur die jeweils
erreichten Programmpfade \cite{nsa-memory-safety}. Auch Zig prüft zur Übersetzungszeit, seine
Laufzeitprüfungen hängen aber vom Build-Modus ab \cite{zig-lang}. Für die Vertrauensgrenze dieser
Bibliothek fiel die Wahl deshalb auf Rust.

Der Parser liest die Optionsbytes als geliehene Slices ohne Kopie; ihre Lebensdauer ist an den
Empfangspuffer gebunden \cite{rust-book}. Zwei
\texttt{unsafe}-Blöcke stehen an den Schnittstellen zum Betriebssystem; ihre Sicherheits\-bedingungen
stehen als Kommentar am jeweiligen Block, ihre Lage im Modul \texttt{socket} nennt \chapref{chap:entwurf}.
Externe Parser-Bibliotheken entfallen, weil diese Arbeit die Mechanik von RFC~9868 selbst untersucht.

Safe Rust verhindert keine Protokollfehler. Ein um ein Byte verschobener Start\-offset bei ungeradem
Beginn der Surplus Area überstand 23 selbstgeschriebene Tests und wurde erst bei der
Pull-Request-Prüfung entdeckt. Der Fehler betraf die Protokolllogik, nicht die Speichersicherheit.
Daraus entstand die Vorgabe, neue oder geänderte Verarbeitungspfade durch Property-Tests und neue oder
erweiterte Fuzz-Ziele abzusichern (\figref{fig:stepzyklus}).

\subsection{Zugang zur Surplus Area}
\label{subsec:surplus-zugang}

Ein gewöhnlicher UDP-Socket unter Linux kann die Surplus Area weder füllen noch auslesen
(\secref{sec:netzpfad}) \cite{linux510udp}. Packet Sockets arbeiten auf Ebene der Netzgeräte,
TUN/TAP benötigt ein virtuelles Netzgerät, und \texttt{AF\_XDP} verlangt für den Empfang ein XDP-Programm
sowie Speicherbereiche und Warteschlangen \cite{man7packet,linux510tuntap,linux510afxdp}.
Ein Kernel\-modul liegt außerhalb des Projektumfangs. Raw Sockets geben der Bibliothek ohne eigenes
Netzgerät Zugriff auf das vollständige IP-Datagramm. Sie sind damit für den gewählten Umfang der
einfachste passende Zugang \cite{man7raw}.

Die Bibliothek nutzt diesen Zugang unter Linux. Sein Verhalten ist dokumentiert und im Kernelquelltext
nachprüfbar; außerdem laufen die Sende- und Empfangs\-programme der Kampagnen unter Linux
(\secref{sec:ff2-pfadmodell}). Im Gegenzug prüft die Bibliothek IPv4-Header, Längen, UDP-Prüfsumme und
Optionen selbst.
Beim Senden liefert sie mit \texttt{IP\_HDRINCL} den IPv4-Header. Da dieser Pfad nicht fragmentiert, muss
jedes gesendete Datagramm mit Optionen in die MTU der ausgehenden Schnittstelle passen (\secref{sec:netzpfad}).
Diese Folgen bestimmen den Entwurf in \chapref{chap:entwurf}.

\subsection{Mitschnitt und Auswertung}
\label{subsec:mitschnittwerkzeuge}

Die Kampagnenskripte zeichnen den Verkehr mit \texttt{tcpdump} in PCAP-Dateien auf. Eigene Python-Programme
werten diese Dateien aus, weil Wireshark in der verwendeten Version~4.6 die UDP-Optionen nicht dekodiert.
Ihr Prüfumfang unterscheidet sich je nach Versuch.

Bei der Wire-Prüfung rechnet \texttt{wire-check.py} die Prüfsummen einschließlich OCS selbst nach und
prüft die Options\-struktur. \texttt{tshark} liefert dazu eine Gegenprobe der IPv4- und UDP-Felder
(\secref{sec:pruefkette}). Auch \texttt{eval-check.py}, das bei den lokalen Referenzläufen und in der
Kampagne vom 10.~August eingesetzt wurde, prüft OCS selbst.

Die Auswerter der späteren Kampagnen, in \secref{sec:testkonzept} als P0, P1 und P2 geführt, verknüpfen
die Sendebelege mit den Paketmitschnitten und den Empfänger\-ausgaben.
Den OCS-Status übernehmen sie aus der Empfänger\-ausgabe; diese Angaben sind deshalb keine
unabhängige OCS-Nachrechnung. Welche Messreihen mit welchen Nachweisen vorliegen, beschreibt
\secref{sec:testkonzept}.

\subsection{Fuzzing und formale Gegenprobe}
\label{subsec:formale-gegenprobe}

Fuzzing und formale Gegenproben ergänzen die Unit-, Integrations- und Property-Tests
(\secref{subsec:pruefmittel}). Sie sollen Fehler aufdecken, die bei der KI-gestützten Entwicklung und
ihren bisherigen Prüfungen unentdeckt bleiben können (\secref{sec:gefahrenmodell}). Fuzz-Ziele und
Property-Tests verwenden dabei gemeinsame Testmodule; fehlerhafte Annahmen in diesen Modulen können deshalb
in beiden Prüfverfahren unentdeckt bleiben.
Das über \texttt{cargo-fuzz} eingebundene libFuzzer verändert Eingaben
des tatsächlichen Parsercodes und bevorzugt solche, die zusätzliche Programmpfade erreichen.
Es sucht nach Abstürzen und verletzten Invarianten \cite{libfuzzer}. Lean beweist ausgewählte Eigenschaften
handgeschriebener Modelle unter den Voraussetzungen der jeweiligen Theoreme; den Begriff führt
\secref{subsec:pruefmittel} ein. Fuzzing beweist keine Fehlerfreiheit, und ein Lean-Beweis gilt nicht
automatisch für den Rust-Code, der das Modell umsetzt. Auch erfolgreiche Prüfungen ersetzen den Abgleich
mit dem RFC nicht. In der Schlussphase deckte ein erneuter Vergleich trotz bestandener Prüfkette weitere
Konformitätslücken auf. Sie wurden in einem eigenen Korrekturschritt bearbeitet (\secref{sec:soll-ist}).
